\section{MPO tensors}

\begin{definition}[MPO tensor]\label{def:mpo_tensor}
    \lean{MPOTensor}
    \leanok
    An \emph{MPO tensor} is a family of matrices
    \[
        \{M^{ij}\}_{i,j=0}^{d-1}, \qquad M^{ij} \in \MN{D},
    \]
    indexed by a ket index $i$ and a bra index $j$. Diagrammatically,
    \begin{center}
        \TNMPOCell{M}{i}{j}
    \end{center}
    The black node denotes the tensor $M$, the horizontal legs are virtual,
    the upper leg is the ket index $i$, and the lower leg is the bra
    index $j$.
\end{definition}

\begin{definition}[Doubled-index MPS view]\label{def:mpo_to_mps}
    \lean{MPOTensor.toMPSTensor}
    \leanok
    \uses{def:mpo_tensor}
    By combining the pair $(i,j)$ into a single physical index in
    $\{0,\ldots,d^2-1\}$, one obtains an MPS tensor with physical dimension
    $d^2$ and bond dimension $D$:
    \[
        M^{\mathrm{MPS}}_{(i,j)} = M^{ij},
    \]
    identifying $\{0,\ldots,d{-}1\} \times \{0,\ldots,d{-}1\}$ with
    $\{0,\ldots,d^2{-}1\}$ via the product encoding.
\end{definition}

\begin{definition}[Word evaluation]\label{def:mpo_eval_word}
    \lean{MPOTensor.evalWord}
    \leanok
    \uses{def:mpo_tensor}
    For words
    $u = (i_1,\ldots,i_N)$ and $v = (j_1,\ldots,j_N)$ in $\{0,\ldots,d{-}1\}^N$,
    the \emph{word evaluation} is
    \[
        M^{u,v} := M^{i_1 j_1} M^{i_2 j_2} \cdots M^{i_N j_N} \in \MN{D}.
    \]
    The empty pair of words evaluates to $\Id_D$, and word pairs of different
    lengths evaluate to $0$.
\end{definition}

\begin{lemma}[Word-evaluation factorization]\label{lem:mpo_eval_word_append}
    \lean{MPOTensor.evalWord_append}
    \leanok
    \uses{def:mpo_eval_word}
    For words $u_1,v_1$ of equal length and arbitrary words $u_2,v_2$,
    word evaluation factors as
    \[
        M^{u_1u_2,\,v_1v_2}=M^{u_1,v_1}M^{u_2,v_2}.
    \]
\end{lemma}

\begin{proof}\leanok\uses{def:mpo_eval_word}
    We argue by induction on the common length of $u_1$ and $v_1$. For the empty prefix,
    \[
        M^{u_2,v_2}=\Id_D\,M^{u_2,v_2}.
    \]
    If $u_1=iu$ and $v_1=jv$, then
    \[
        M^{iuu_2,\,jvv_2}
        = M^{ij}M^{uu_2,\,vv_2}
        = M^{ij}M^{u,v}M^{u_2,v_2}
        = M^{iu,jv}M^{u_2,v_2},
    \]
    using associativity of matrix multiplication.
\end{proof}

\begin{lemma}[Trace cyclicity of the closed MPO word]\label{lem:mpo_trace_evalWord_cyclic}
    \lean{MPOTensor.trace_evalWord_cons_eq_append}
    \leanok
    \uses{def:mpo_eval_word}
    Let $a,b\in\{0,\ldots,d{-}1\}$, and let $u,v$ be words of equal length.
    Writing $au,bv$ for adding the letters at the beginning and $ua,vb$ for
    adding them at the end,
    \[
        \tr\!\left(M^{au,\,bv}\right)
        =
        \tr\!\left(M^{ua,\,vb}\right).
    \]
\end{lemma}

\begin{proof}\leanok\uses{def:mpo_eval_word, lem:mpo_eval_word_append}
    The factorization Lemma~\ref{lem:mpo_eval_word_append} gives
    \[
        M^{ua,\,vb}=M^{u,v}M^{ab},
        \qquad
        M^{au,\,bv}=M^{ab}M^{u,v}.
    \]
    Hence
    \[
        \tr\!\left(M^{ua,\,vb}\right)
        =\tr\!\left(M^{u,v}M^{ab}\right)
        =\tr\!\left(M^{ab}M^{u,v}\right)
        =\tr\!\left(M^{au,\,bv}\right),
    \]
    by cyclicity of the trace.
\end{proof}

\begin{definition}[MPO operator family]\label{def:mpo_family}
    \lean{MPOTensor.mpo}
    \leanok
    \uses{def:mpo_eval_word}
    The operator generated by $M$ on $N$ sites is the matrix
    $\rho^{(N)}(M)$ indexed by configurations
    $\sigma,\tau \in \{0,\ldots,d{-}1\}^N$ and given by
    \[
        \rho^{(N)}(M)_{\sigma,\tau} = \tr(M^{\sigma,\tau}).
    \]
\end{definition}

Diagrammatically,
\begin{center}
    \TNMPOChain{M}{\sigma_1}{\tau_1}{\sigma_N}{\tau_N}{N}
\end{center}
each black node denotes the same local tensor $M$, and the matrix entry
$\rho^{(N)}(M)_{\sigma,\tau}$ is obtained by closing the outer virtual
legs cyclically and taking the resulting trace.

\begin{definition}[Cyclic-shift reindexing]\label{def:mpo_rotate_config}
    \lean{MPOTensor.rotateConfig}
    \leanok
    The \emph{cyclic shift} of configurations on $N$ sites is the bijection that
    relabels each site by one position, $\sigma \mapsto \sigma \circ r$, where
    $r$ is the cyclic rotation $i \mapsto i + 1$ of $\{0,\ldots,N-1\}$.
\end{definition}

\begin{theorem}[Translation invariance of the periodic MPDO]
    \label{thm:mpo_translation_invariant}
    \lean{MPOTensor.mpo_submatrix_rotateConfig}
    \leanok
    \uses{def:mpo_family, def:mpo_rotate_config}
    The operator $\rho^{(N)}(M)$ is invariant under the simultaneous cyclic shift
    of its bra and ket configurations:
    \[
        \rho^{(N)}(M)_{r\sigma,\, r\tau} = \rho^{(N)}(M)_{\sigma,\tau}.
    \]
\end{theorem}

\begin{proof}\leanok
    The entry is a closed bond trace,
    $\rho^{(N)}(M)_{\sigma,\tau}
        = \tr\bigl(M^{\sigma_1\tau_1} \cdots M^{\sigma_N\tau_N}\bigr)$.
    Shifting both configurations by $r$ rotates the matrix product by one factor,
    and the trace is invariant under that cyclic rotation,
    \[
        \tr\bigl(M^{\sigma_2\tau_2} \cdots M^{\sigma_N\tau_N} M^{\sigma_1\tau_1}\bigr)
        = \tr\bigl(M^{\sigma_1\tau_1} \cdots M^{\sigma_N\tau_N}\bigr),
    \]
    so $\rho^{(N)}(M)_{r\sigma,\,r\tau} = \rho^{(N)}(M)_{\sigma,\tau}$.
\end{proof}

\begin{definition}[Hermitian MPO tensor]\label{def:mpo_hermitian}
    \lean{MPOTensor.IsHermitian}
    \leanok
    \uses{def:mpo_tensor}
    An MPO tensor is \emph{Hermitian} if, for all
    $i,j \in \{0,\ldots,d{-}1\}$,
    \[
        M^{ij} = (M^{ji})^\dagger
    \]
\end{definition}

\section{Transfer maps}

\begin{definition}[MPO transfer map]\label{def:mpo_transfer_map}
    \lean{MPOTensor.transferMap}
    \leanok
    \uses{def:mpo_tensor}
    The \emph{transfer map} associated to $M$ is the linear map
    $\E_M : \MN{D} \to \MN{D}$ defined by
    \[
        \E_M(X) = \sum_{i,j=0}^{d-1} M^{ij} X (M^{ij})^\dagger.
    \]
\end{definition}

\begin{lemma}[Identification with the doubled-index MPS transfer map]\label{lem:mpo_transfer_eq_mps}
    \lean{MPOTensor.transferMap_eq_toMPSTensor}
    \leanok
    \uses{def:mpo_to_mps, def:mpo_transfer_map, def:transfer_map}
    The transfer map of an MPO tensor agrees with the transfer map of its
    doubled-index MPS view.
\end{lemma}

\begin{proof}
    \leanok
    \uses{def:mpo_to_mps, def:mpo_transfer_map, def:transfer_map}
    After identifying the single physical index of the doubled tensor with a
    pair $(i,j)$, the single sum defining the MPS transfer map is exactly the
    double sum defining $\E_M$.
\end{proof}

\begin{theorem}[Positivity of the MPO transfer map]\label{thm:mpo_transfer_pos}
    \lean{MPOTensor.transferMap_pos}
    \leanok
    \uses{def:mpo_transfer_map}
    If $X \ge 0$, then $\E_M(X) \ge 0$.
\end{theorem}

\begin{proof}
    \leanok
    \uses{lem:mpo_transfer_eq_mps}
    By Lemma~\ref{lem:mpo_transfer_eq_mps}, $\E_M$ is the transfer map of the
    doubled-index MPS tensor. The corresponding MPS transfer map is positive,
    so $\E_M(X) \ge 0$ whenever $X \ge 0$.
\end{proof}

\section{Zero correlation length}

\begin{definition}[Zero correlation length for MPO tensors]%
    \label{def:mpo_is_zcl}
    \lean{MPOTensor.IsZCL}
    \leanok
    \uses{def:mpo_transfer_map}
    An MPO tensor $M$ has \emph{zero correlation length} when its transfer map
    is idempotent:
    \[
        \E_M \circ \E_M = \E_M.
    \]
    This is the mixed-state analog of the renormalization fixed-point
    condition for MPS tensors; see \cite[definition of MPDO zero correlation
        length, lines~736--741]{Cirac2016MPDO_arXiv} and
    \cite[Section~II.E.2, lines~937--939]{Cirac2021Matrix}.
\end{definition}

\begin{theorem}[MPO ZCL iff doubled-index MPS RFP]%
    \label{thm:mpo_is_zcl_iff_to_mps_rfp}
    \lean{MPOTensor.isZCL_iff_toMPSTensor_isRFP}
    \leanok
    \uses{def:mpo_is_zcl, def:mpo_to_mps, def:is_rfp, lem:mpo_transfer_eq_mps}
    An MPO tensor has zero correlation length if and only if its doubled-index
    MPS view is a renormalization fixed point.
\end{theorem}

\begin{proof}
    \leanok
    By Lemma~\ref{lem:mpo_transfer_eq_mps}, the transfer map of $M$ agrees with
    the transfer map of its doubled-index MPS view. Thus
    $\E_M \circ \E_M = \E_M$ holds exactly when the transfer map of
    that doubled-index MPS tensor is idempotent, which is the RFP condition.
\end{proof}

\section{MPDOs and LPDOs}

\begin{definition}[MPDO]\label{def:mpdo}
    \lean{MPOTensor.IsMPDO}
    \leanok
    \uses{def:mpo_family}
    An MPO tensor is an \emph{MPDO} if, for every $N \in \N$,
    \[
        \rho^{(N)}(M) \ge 0
    \]
    This is the global positivity condition of \cite[Section~4]{Cirac2016MPDO_arXiv}.
\end{definition}

\begin{definition}[LPDO]\label{def:lpdo}
    \lean{MPOTensor.IsLPDO}
    \leanok
    \uses{def:mpo_tensor}
    An MPO tensor is an \emph{LPDO} if there exist a Kraus dimension $d_K$,
    an inner bond dimension $D'$, a bond-space identification
    $e : \{0, \ldots, D - 1\} \simeq \{0, \ldots, D' - 1\} \times \{0, \ldots, D' - 1\}$
    (equivalently, $D$ and $(D')^2$
    have the same cardinality), and matrices $A^{(i,k)} \in \MN{D'}$ such that,
    for all physical indices $i,j$,
    \[
        M^{ij}
        = \left(\sum_{k=0}^{d_K-1} A^{(i,k)} \otimes \overline{A^{(j,k)}}\right)_{e,e}
    \]
    where $\otimes$ denotes the Kronecker product, $\overline{(\cdot)}$
    denotes entrywise complex conjugation, and $(\cdot)_{e,e}$ denotes
    reindexing along $e$ back to $\{0, \ldots, D - 1\}$.
    See \cite[Section~4.3]{Cirac2016MPDO_arXiv}.
\end{definition}

\begin{theorem}[LPDOs are MPDOs]\label{thm:lpdo_is_mpdo}
    \lean{MPOTensor.IsLPDO.isMPDO}
    \leanok
    \uses{def:lpdo, def:mpdo}
    Every LPDO tensor is an MPDO.
\end{theorem}

\begin{proof}
    \leanok
    \uses{def:lpdo, def:mpdo}
    By the mixed-product property
    $(A \otimes B)(C \otimes D) = (AC) \otimes (BD)$, the $N$-site product
    decomposes as a sum of Kronecker products after transporting along the
    bond-space identification $e$:
    \[
        M^{\sigma_1 \tau_1} \cdots M^{\sigma_N \tau_N}
        = (\sum_{\kappa} P_\kappa \otimes \overline{Q_\kappa})
        \big|_{e},
    \]
    where $P_\kappa = A^{(\sigma_1,\kappa_1)} \cdots A^{(\sigma_N,\kappa_N)}$
    and $Q_\kappa = A^{(\tau_1,\kappa_1)} \cdots A^{(\tau_N,\kappa_N)}$.
    Since the trace is invariant under the reindexing by~$e$,
    we apply $\tr(X \otimes Y) = \tr(X)\tr(Y)$ to obtain
    \[
        \rho^{(N)}(M)_{\sigma,\tau}
        = \sum_\kappa \tr(P_\kappa) \overline{\tr(Q_\kappa)}
        = \sum_\kappa \psi_\kappa(\sigma) \overline{\psi_\kappa(\tau)},
    \]
    where $\psi_\kappa(\sigma) = \tr(P_\kappa)$.
    Hence $\rho^{(N)}(M) = \sum_\kappa \ketbra{\psi_\kappa}{\psi_\kappa} \ge 0$.
\end{proof}

\begin{definition}[Local purification RFP condition]\label{def:is_local_purification_rfp}
    \lean{MPOTensor.IsLocalPurificationRFP}
    \leanok
    \uses{def:lpdo, def:is_rfp}
    An MPO tensor $M$ satisfies the local purification RFP condition if it is an
    LPDO whose purifying tensor $A$, viewed as a matrix product state tensor on
    the combined spin--ancilla index, is a pure-state renormalization fixed point.
    Thus
    \[
        M^{ij}
        =
        \left(\sum_{k} A^{(i,k)} \otimes \overline{A^{(j,k)}}\right)_{e,e}.
    \]
    This local condition is motivated by the purification tensor formula
    \cite[lines~744--747]{Cirac2016MPDO_arXiv}, but is weaker than the source
    purification RFP definition; the latter remains open.
\end{definition}

\begin{theorem}[A local purification RFP generates MPDOs]
    \label{thm:local_purification_rfp_is_mpdo}
    \lean{MPOTensor.IsLocalPurificationRFP.isMPDO}
    \leanok
    \uses{def:is_local_purification_rfp, def:mpdo, thm:lpdo_is_mpdo}
    Every tensor satisfying the local purification RFP condition generates
    matrix product density operators.
\end{theorem}

\begin{proof}\leanok
    \uses{thm:lpdo_is_mpdo}
    Dropping the renormalization-fixed-point condition on the purifying tensor
    leaves the local purification structure, so $M$ is an LPDO; the conclusion
    follows from Theorem~\ref{thm:lpdo_is_mpdo}.
\end{proof}

\begin{theorem}[Local purification RFP does not imply zero correlation length]
    \label{thm:local_purification_rfp_not_zcl}
    \lean{MPOTensor.exists_isLocalPurificationRFP_not_isZCL}
    \leanok
    \uses{def:is_local_purification_rfp, def:mpo_is_zcl}
    There is an MPO tensor satisfying the local purification RFP condition for
    which the literal transfer-map idempotence $\E_M\circ\E_M=\E_M$ fails.
\end{theorem}

% Counterexample analysed in docs/paper-gaps/cpsv16_zcl_canonical_form_normalization.tex.
\begin{proof}\leanok
    Take the diagonal purification at $d=d_K=2$ and $D=D'=1$ with amplitudes
    $A=[\tfrac{1}{\sqrt2},0,0,\tfrac{1}{\sqrt2}]$. The purifying tensor is a
    pure-state renormalization fixed point, since $\sum|A|^2=1$ makes its transfer
    map the identity. The ancilla contraction
    \[
        M^{ij}=\sum_k A^{(i,k)}\,\overline{A^{(j,k)}}
    \]
    gives the scalar entries $M^{00}=M^{11}=\tfrac12$ (off-diagonal $0$), so the
    induced transfer map is $\tfrac12\cdot\mathrm{id}$ and
    $\E_M\circ\E_M=\tfrac14\cdot\mathrm{id}\neq\E_M$. The trace contraction has
    dropped the leading eigenvalue from $1$ to $\tfrac12$; the literal
    zero-correlation-length condition is therefore strictly stronger than the
    source's normalized one.
\end{proof}

